% Source: https://artofproblemsolving.com/community/c3h1594200_nondegenerate_triangles_from_a_set
% Stable identity: aops-topic-1594200-q01
Let the sides be $a<b<c$ from $\{1,2,3,4,5\}$. The triangle is non-degenerate exactly when $a+b>c$ (the other two triangle inequalities hold automatically, since $c$ is the largest side). Go through the largest side $c$:
\begin{itemize}
\item $c=3$: $(a,b)=(1,2)$, but $1+2=3$ is not greater than $3$ (degenerate). None.
\item $c=4$: $(1,2),(1,3),(2,3)$; only $2+3>4$ works. One: $(2,3,4)$.
\item $c=5$: $(a,b)$ from $\{1,2,3,4\}$ with $a+b>5$: $(2,4)$ and $(3,4)$. Two: $(2,4,5)$, $(3,4,5)$.
\end{itemize}
So there are $\boxed{3}$ such triangles:

\begin{asy}
usepackage("amsmath");
size(8cm);
defaultpen(fontsize(10pt));
pair tri(real a,real b,real c){ return (( a*a+b*b-c*c)/(2*a), sqrt(b*b-((a*a+b*b-c*c)/(2*a))^2)); }
void T(real a,real b,real c,pair o){
  pair P=o,Q=o+(a,0),R=o+tri(a,b,c);
  draw(P--Q--R--cycle,linewidth(0.9));
  pair n1=unit(rotate(-90)*(Q-P)),n2=unit(rotate(-90)*(R-Q)),n3=unit(rotate(-90)*(P-R));
  label("$"+string(a)+"$",(P+Q)/2,n1);label("$"+string(c)+"$",(Q+R)/2,n2);label("$"+string(b)+"$",(R+P)/2,n3);
}
T(4,2,3,(0,0));T(5,2,4,(6,0));T(5,3,4,(13,0));
label("\mbox{}",(min(currentpicture,user=true)+max(currentpicture,user=true))/2,black); // keeps the text after the figure black when LaTeX typesets the labels inline
\end{asy}
